\chapter{Geometry of the dipole flow}
\label{chap:geometry}

This chapter collects explicit facts about $\Psi=\cos x+\cos y$ and its flow.
Most are elementary trigonometric identities or one-variable estimates and are
good early formalization targets.

\begin{remark}[Digest note: notation]
The manuscript reuses letters heavily across sections. The level variable is
$r=\Psi$ in the even--even section, $s=\Psi$ in the odd--odd section (where
$r=|\nabla\Psi|^2$ instead), and $r$ denotes distance to a saddle in the
regularity section. Similarly $h$ is $\Psi/2$, $\cos a$, or a generic level,
depending on the section. In this blueprint we write $\Psi$ for the stream
function, $h=\Psi/2$ for the level on the positive cell, $\rho=|\nabla\Psi|$,
and $d$ for distance to the saddle set. Local deviations are flagged.
\end{remark}

\section{Cells, separatrix, critical points}

\begin{definition}[Cells and diagonal coordinates]
\label{def:cells}
\uses{def:dipole}
The critical points of $\Psi$ are the elliptic centers $(0,0)$ (maximum,
$\Psi=2$) and $(\pi,\pi)$ (minimum, $\Psi=-2$), and the two saddles $(\pi,0)$,
$(0,\pi)$ (where $\Psi=0$). The separatrix $\{\Psi=0\}$ is the union of the
two closed lines $\{x+y\equiv\pi\}$ and $\{x-y\equiv\pi\}$. The positive cell
is $\Cell=\{\Psi>0\}$, the negative cell $\{\Psi<0\}$.

On a lift of $\Cell$ put $a=(x+y)/2$, $b=(x-y)/2$. Then $\Cell$ is the square
$|a|,|b|<\pi/2$, and $\Psi=2\cos a\cos b$, $dx\,dy=2\,da\,db$, and
$\Delta=\tfrac12(\partial_a^2+\partial_b^2)$. The edges $a=\pm\pi/2$ lie on the
fixed line of $S_1$ and the edges $b=\pm\pi/2$ on the fixed line of $S_2$.
The four corners $(a,b)=(\pm\pi/2,\pm\pi/2)$ project to the two saddles, each
saddle being a corner twice.
\end{definition}

\begin{definition}[Classical operators on the lift]
\label{def:lift-operators}
\uses{def:dipole}
\lean{NoEmbeddedEigenvalues.partialX,
  NoEmbeddedEigenvalues.partialY,
  NoEmbeddedEigenvalues.laplacian,
  NoEmbeddedEigenvalues.classicalL0,
  NoEmbeddedEigenvalues.gradDot,
  NoEmbeddedEigenvalues.gradSq,
  NoEmbeddedEigenvalues.dipoleLift}
\leanok
Pointwise differential identities are stated for functions on the real lift
$\R^2$ of $\T^2$: $\partial_x,\partial_y$ are ordinary partial derivatives,
$\Delta=\partial_x^2+\partial_y^2$, $L_0=\sin y\,\partial_x-\sin x\,\partial_y$,
$\nabla u\cdot\nabla v=u_xv_x+u_yv_y$ and $|\nabla u|^2=\nabla u\cdot\nabla u$.
The dipole lifts to $\Psi(x,y)=\cos x+\cos y$ on $\R^2$.
\end{definition}

\begin{lemma}[Calculus on the lift]
\label{lem:lift-calculus}
\uses{def:lift-operators}
\lean{NoEmbeddedEigenvalues.partialX_add,
  NoEmbeddedEigenvalues.partialX_sub,
  NoEmbeddedEigenvalues.partialX_smul,
  NoEmbeddedEigenvalues.partialX_const_smul,
  NoEmbeddedEigenvalues.partialY_add,
  NoEmbeddedEigenvalues.partialY_sub,
  NoEmbeddedEigenvalues.partialY_smul,
  NoEmbeddedEigenvalues.partialY_const_smul,
  NoEmbeddedEigenvalues.partialX_eq_zero_of_eventually,
  NoEmbeddedEigenvalues.partialY_eq_zero_of_eventually,
  NoEmbeddedEigenvalues.partialY_partialX,
  NoEmbeddedEigenvalues.contDiff_partialX,
  NoEmbeddedEigenvalues.contDiff_partialY}
\leanok
The partial derivatives satisfy the sum and product rules at points of
differentiability, vanish at $p$ for functions vanishing near $p$, map $C^{n+1}$
to $C^n$, and commute on $C^2$ functions: $\partial_y\partial_xu=\partial_x\partial_yu$.
\end{lemma}

\begin{proof}
\leanok
One-variable calculus along the coordinate lines. The partials are the
Fréchet derivative applied to $(1,0)$ and $(0,1)$, so symmetry of mixed
partials is symmetry of the second Fréchet derivative of a $C^2$ function.
\end{proof}

\begin{definition}[Integration over the torus on the lift]
\label{def:square-integral}
\uses{def:lift-operators}
\lean{NoEmbeddedEigenvalues.IsPeriodicLift,
  NoEmbeddedEigenvalues.periodSquare,
  NoEmbeddedEigenvalues.squareIntegral,
  NoEmbeddedEigenvalues.squarePairing}
\leanok
A function on $\R^2$ is a \emph{periodic lift} if it is $2\pi$-periodic in both
variables. For such functions, $\int_{\T^2}g$ means $\int_{[0,2\pi]^2}g$, and
$\ip uv=\int_{[0,2\pi]^2}u\bar v$.
\end{definition}

\begin{lemma}[Integrals of derivatives vanish]
\label{lem:square-integral-partial}
\uses{def:square-integral, lem:lift-calculus}
\lean{NoEmbeddedEigenvalues.squareIntegral_partialX,
  NoEmbeddedEigenvalues.squareIntegral_partialY,
  NoEmbeddedEigenvalues.IsPeriodicLift.partialX,
  NoEmbeddedEigenvalues.IsPeriodicLift.partialY}
\leanok
Partial derivatives of periodic lifts are periodic lifts, and
$\int_{\T^2}\partial_xg=\int_{\T^2}\partial_yg=0$ for every periodic $C^1$ lift $g$.
\end{lemma}

\begin{proof}
\leanok
Fubini, and on each coordinate line the fundamental theorem of calculus:
$\int_0^{2\pi}\partial_yg(x,y)\,dy=g(x,2\pi)-g(x,0)=0$.
\end{proof}

\begin{lemma}[The pairing]
\label{lem:square-pairing}
\uses{def:square-integral}
\lean{NoEmbeddedEigenvalues.squarePairing_add_left,
  NoEmbeddedEigenvalues.squarePairing_sub_left,
  NoEmbeddedEigenvalues.squarePairing_const_mul_left,
  NoEmbeddedEigenvalues.squarePairing_const_mul_right,
  NoEmbeddedEigenvalues.squarePairing_comm,
  NoEmbeddedEigenvalues.squarePairing_real_smul}
\leanok
For continuous periodic lifts, $\ip uv$ is additive and $\C$-linear in $u$,
conjugate linear in $v$, $\ip vu=\overline{\ip uv}$, and $\ip{fu}v=\ip u{fv}$
for real $f$.
\end{lemma}

\begin{proof}
\leanok
Continuous functions are integrable on the compact square.
\end{proof}

\begin{lemma}[Algebraic identities for the dipole]
\label{lem:dipole-identities}
\uses{def:dipole, def:lift-operators}
\lean{NoEmbeddedEigenvalues.Ψ_coe,
  NoEmbeddedEigenvalues.partialX_dipoleLift,
  NoEmbeddedEigenvalues.partialY_dipoleLift,
  NoEmbeddedEigenvalues.gradSq_dipoleLift,
  NoEmbeddedEigenvalues.laplacian_dipoleLift,
  NoEmbeddedEigenvalues.classicalL0_dipoleLift,
  NoEmbeddedEigenvalues.gradSq_dipoleLift_sub,
  NoEmbeddedEigenvalues.dipoleLift_mul_le_gradSq,
  NoEmbeddedEigenvalues.sq_cos_sub_cos,
  NoEmbeddedEigenvalues.gradDot_dipoleLift_gradSq,
  NoEmbeddedEigenvalues.laplacian_gradSq_dipoleLift,
  NoEmbeddedEigenvalues.gradSq_gradSq_dipoleLift,
  NoEmbeddedEigenvalues.laplacian_sin_fst,
  NoEmbeddedEigenvalues.laplacian_sin_snd}
\leanok
Write $s=\Psi$, $\rho^2=|\nabla\Psi|^2$, $\delta=\cos x-\cos y$. Then
\begin{align*}
  &\Delta\Psi=-\Psi,\qquad L_0\Psi=0,\qquad \rho^2=\sin^2x+\sin^2y,\\
  &\rho^2-s(2-s)=2(1-\cos x)(1-\cos y)\ \ge0,\qquad
   \delta^2=4-2\rho^2-s^2,\\
  &\nabla s\cdot\nabla\rho^2=s(4-s^2-3\rho^2),\qquad
   \Delta\rho^2=4(1-\rho^2),\\
  &|\nabla\rho^2|^2=4\rho^2-2\rho^4+4\rho^2s^2+2s^4-8s^2.
\end{align*}
The coefficients of $L_0$, namely $\sin y$ and $\sin x$, are first Laplace
eigenfunctions: $\Delta\sin x=-\sin x$.
(Manuscript \texttt{(ee:c-positive)} and the list before the coefficient
table in Section~5.)
\end{lemma}

\begin{proof}
\leanok
Trigonometric identities. For example, with $p=\cos x$ and $q=\cos y$,
$\rho^2=2-p^2-q^2$ and $\nabla s\cdot\nabla\rho^2=-2p(1-p^2)-2q(1-q^2)$.
\end{proof}

\section{The separatrix arcs}

\begin{definition}[Heteroclinic arc]
\label{def:separatrix-arc}
\uses{def:cells}
\lean{NoEmbeddedEigenvalues.arcAngle,
  NoEmbeddedEigenvalues.separatrixArc}
\leanok
Let $\gamma(s)=\big(2\arctan(e^{-s}),\,2\arctan(e^{-s})-\pi\big)$, $s\in\R$
(manuscript \texttt{(ee:arc)}).
\end{definition}

\begin{remark}
The images of $\gamma$ under the rotation $R$ give the four directed arcs
bounding $\Cell$. This is not yet a separate node or formalized.
\end{remark}

\begin{lemma}[Arc parametrization]
\label{lem:separatrix-arc}
\uses{def:separatrix-arc}
\lean{NoEmbeddedEigenvalues.sin_arcAngle,
  NoEmbeddedEigenvalues.hasDerivAt_arcAngle,
  NoEmbeddedEigenvalues.separatrixArc_velocity,
  NoEmbeddedEigenvalues.hasDerivAt_separatrixArc,
  NoEmbeddedEigenvalues.separatrixArc_speed,
  NoEmbeddedEigenvalues.separatrixArc_fst_sub_snd,
  NoEmbeddedEigenvalues.tendsto_separatrixArc_atBot,
  NoEmbeddedEigenvalues.tendsto_separatrixArc_atTop,
  NoEmbeddedEigenvalues.hasDerivAt_tanh,
  NoEmbeddedEigenvalues.one_sub_tanh_sq}
\leanok
$\gamma$ is an integral curve of the velocity $(\sin y,-\sin x)$, lies on
$\{x-y=\pi\}$, runs from the saddle $(\pi,0)$ ($s\to-\infty$) to the saddle
$(0,-\pi)\equiv(0,\pi)$ ($s\to+\infty$), and has speed
$|\gamma'(s)|=\sqrt2\sech s$ (manuscript \texttt{(ee:arc-speed)}). In the
variable $\xi=\tanh s$ one has $ds=d\xi/(1-\xi^2)$.
\end{lemma}

\begin{proof}
\leanok
With $x=2\arctan(e^{-s})$ one has $\sin x=\sech s$ and $x'=-\sech s$; since
$y=x-\pi$, $\sin y=-\sin x$.
\end{proof}

\section{Orbits near the separatrix}

\begin{definition}[Quarter orbits and the period]
\label{def:period}
\uses{def:cells}
For $0<h<1$ let $\gamma_h$ be the orbit on the level $\{\Psi=2h\}$ in $\Cell$,
parametrized by flow time and centered at $(a,b)=(0,\arccos h)$, i.e. at
$(x,y)=(\arccos h,-\arccos h)$. One quarter of the orbit is
$\gamma_h([-L(h),L(h)])$, where
\[
  L(h)=\int_0^{\arccos\sqrt h}\frac{da}{\sqrt{\cos^2a-h^2}}
      =\int_0^1\frac{dt}{\sqrt{1-(1-h)t^2}\sqrt{1+h-t^2}}
\]
(manuscript \texttt{(ee:period)}). The full period is $S(h)=8L(h)$, and the
orbital frequency is $b(h)=2\pi/S(h)$.
\end{definition}

\begin{lemma}[The period function]
\label{lem:period}
\uses{def:period}
$L$ is strictly decreasing on $(0,1)$, $S(h)=4\mathbf K(1-h^2)$ with
$\mathbf K(m)=\int_0^{\pi/2}(1-m\sin^2\theta)^{-1/2}d\theta$, and as $h\to0$,
\[
  L(h)=\tfrac12\log(4/h)+O(h^2\log(1/h)),\qquad
  L'(h)=-\frac1{2h}+O(h\log(1/h)).
\]
In particular $e^{L(h)}\le Ch^{-1/2}$, and $S$ has simple level sets.
\end{lemma}

\begin{proof}
Monotonicity: the derivative in $h$ of the squared denominator is
$t^2(1+h-t^2)+1-t^2+ht^2>0$. The asymptotics follow by substituting $y=1-t^2$
and splitting at $y=h$; the precise expansion follows from the classical
expansion of $\mathbf K$ near $m=1$.
\end{proof}

\begin{lemma}[Orbit comparison]
\label{lem:orbit-comparison}
\uses{def:period, def:separatrix-arc}
Let $d$ denote distance to the saddle set. For $|t|\le L(h)$ and small $h$:
\begin{enumerate}
  \item $|\gamma_h(t)-\gamma(t)|\le C\sqrt h$ (after matching the centering of
        $\gamma$ and $\gamma_h$) (manuscript \texttt{(ee:orbit-comparison)});
  \item $ce^{-|t|}\le d(\gamma_h(t))\le Ce^{-|t|}$; in particular both
        endpoints of the quarter orbit are at distance $\sim\sqrt h$ from a
        saddle;
  \item $|\partial_h\gamma_h(t)|\le Ce^{|t|}\le C/d$ and
        $|\partial_h^2\gamma_h(t)|\le Ce^{2|t|}\le C/d^2$
        (manuscript \texttt{(eq:quarter-orbit-bounds)});
  \item in the orbit coordinates $(h,t)$ the area element is $2\,dh\,dt$, and
        $|\nabla h|\le Cd$, $|\nabla t|\le C/d$;
  \item $\int_{-L}^{L}d^{-2}\,dt\le C/h$, $\int_{-L}^{L}d^{-4}\,dt\le C/h^2$,
        and $\int_{-L}^{L}d^{\,p}\,dt\le C$ for $p>0$.
\end{enumerate}
\end{lemma}

\begin{proof}
\uses{lem:period}
The velocity field has Lipschitz constant at most one; the midpoint
displacement is $O(h)$ and grows at most like $e^{|t|}$, while $he^{2L}$ is
bounded. The remaining bounds follow from comparison with the hyperbolic
linearization at the saddles.
\end{proof}

\begin{remark}[Digest note]
Items (3)--(5) are used throughout Chapter~\ref{chap:regularity} but are only
asserted in the manuscript. A cleaner route may be the exact corner coordinates
of Lemma~\ref{lem:corner-coordinates}, in which level sets are exact
hyperbolas.
\end{remark}

\begin{lemma}[Corner coordinates]
\label{lem:corner-coordinates}
\uses{def:cells}
Near a corner of $\Cell$, let $\alpha,\beta>0$ be the distances in the $(a,b)$
coordinates to the two sides meeting there, and put $u=\sin\alpha$,
$v=\sin\beta$. Then $\Psi/2=uv$ exactly, the metric coefficients in $(u,v)$
are smooth, even in each variable, and uniformly comparable to constants on a
fixed small corner. Along the level $uv=h$, flow-time measure $dt$ is
comparable to $du/u$.
\end{lemma}

\begin{proof}
$\cos a=\sin\alpha$ and $\cos b=\sin\beta$ at the corner $(\pi/2,\pi/2)$, and
similarly at the other corners.
\end{proof}

\begin{lemma}[Saddle normal form]
\label{lem:saddle-normal-form}
\uses{def:cells}
Near each saddle there are orthonormal coordinates $(u,v)$, in which the two
separatrix rays are the coordinate axes, such that
$\Psi=uv+O(r^4)$ and $|\nabla\Psi|^2=u^2+v^2+O(r^4)$, $r=(u^2+v^2)^{1/2}$.
The linearization $H_s$ of $H$ at the saddle has eigenvalues $\pm i$ on linear
polynomials.
\end{lemma}

\begin{proof}
Taylor expansion of $\cos x+\cos y$ at $(\pi,0)$ in the rotated coordinates.
\end{proof}
